\documentclass[14pt]{article}
\usepackage{graphicx} % Required for inserting images
\usepackage{titlesec}
\usepackage{multicol, caption}
\usepackage{blindtext}
\usepackage{amsmath}
\usepackage{titlesec}
\usepackage{gensymb}

\usepackage{amsmath}
\usepackage{amssymb}
\usepackage{cite}
\usepackage{caption}
\usepackage[a4paper,total={6.5in,10in}]{geometry}
\usepackage{amssymb}
\usepackage{natbib}
\usepackage[colorlinks=true, allcolors=black]{hyperref}
\usepackage{titlesec}
\usepackage{titling}
\usepackage{wrapfig}
\usepackage{lipsum}
\usepackage{fancyhdr}
\bibliographystyle{IEEEtran}
\setcounter{tocdepth}{2}
\titleformat*{\subsubsection}{\normalfont\fontfamily{cmr}\fontsize{12}{12}\selectfont}

\setlength{\droptitle}{-5em} 
\graphicspath{ {./Images/} }
\setlength\columnsep{30pt}
\newenvironment{Figure}
  {\par\medskip\noindent\minipage{\linewidth}}
  {\endminipage\par\medskip}

\titleformat*{\section}{%
    \fontsize{16}{17}\bfseries%
}

\titleformat*{\subsection}{%
    \fontsize{13}{15}\bfseries%
}
\title{CIV102F - Structures and Materials Equations}
\author{Miles Ehrlich}
\date{December 2025}

\begin{document}

\maketitle
\begin{multicols}{2}

    \section{Basic Concepts}
    \subsection{Equilibrium}
    
    \begin{equation}
        \\ \sum F_x=0 \quad
        \\ \sum F_y=0 \quad
        \\ \sum M_P=0
    \end{equation}
    \begin{equation}
        F_y=F_x \text{tan}(\theta)
    \end{equation}

    \subsection{Moment}
    \begin{equation}
        M=F\cdot d_{\perp}
    \end{equation}
    \subsection{Stress-Strain}
    \begin{equation}
        \sigma = \frac{F}{A}
    \end{equation}
    \begin{equation}
        \varepsilon=\frac{\Delta}{L_0}
    \end{equation}
    \begin{equation}
        \sigma=E\cdot \varepsilon=k\frac{L_0}{A}\cdot \varepsilon
    \end{equation}
    \begin{equation}
        \Delta=\frac{FL_0}{EA}
    \end{equation}
    \begin{equation}
        U=\int \sigma d\varepsilon=\frac{1}{2}\sigma \varepsilon
    \end{equation}
    \begin{equation}
        W=U\cdot V_0
    \end{equation}
    \begin{equation}
        R=\int \limits_{0}^{\varepsilon_y} \sigma d\varepsilon
    \end{equation}
    \begin{equation}
        T=\int \limits_{0}^{\varepsilon_r} \sigma d\varepsilon
    \end{equation}
    \section{Suspension Bridges}
    \begin{equation}
        A_y=B_y=\frac{wL}{2}
    \end{equation}
    \begin{equation}
        H=\frac{wL^2}{8h}
    \end{equation}
    \begin{equation}
        T_{max}=\sqrt{H^2+A_y^2}
    \end{equation}
    \section{Structural Analysis}
    \subsection{Pulleys}
    \begin{equation}
        T1=T2
    \end{equation}
    \subsection{Rollers}
    \begin{equation}
        \sum F_y=0
    \end{equation}
    \subsection{Pins, Hinges}
    \begin{equation}
        \sum F_x = \sum F_y = 0
    \end{equation}
    \subsection{Fixed End}
    \begin{equation}
        \sum F_x = \sum F_y = \sum M=0
    \end{equation}
    \section{Oscillations}
    \subsection{Simple Harmonic Motion}
    \begin{equation}
        F=-kx
    \end{equation}
    \begin{equation}
        x(t)=A\text{sin}(\omega t) + \Delta_0
    \end{equation}
    \begin{equation}
        \omega=\sqrt{\frac{k}{m}}
    \end{equation}
    \begin{equation}
        f_{n,point}=\frac{15.76}{\sqrt{\Delta_0}}
    \end{equation}
    \begin{equation}
        f_{n,UDL}=\frac{17.76}{\sqrt{\Delta_0}}
    \end{equation}
    \subsection{Driven Damped Oscillation}
    \begin{equation}
        F(t)=F_0\text{sin}(\omega t)
    \end{equation}
    \begin{equation}
        x(t)=DAF\cdot \frac{F_0}{k} \text{sin}(\omega t + \phi)+\Delta_0
    \end{equation}
    \begin{equation}
        DAF=\frac{1}{\sqrt{\left(1-\left(\frac{f}{f_n}\right)^2\right)^2+\left(\frac{2\beta f}{f_n}\right)^2}}
    \end{equation}
    \begin{equation}
        F_{max}=F_{stat}+DAF\cdot F_0
    \end{equation}
    \begin{equation}
        \Delta x_{max} = \Delta x_{stat} \cdot \frac{F_{max}}{F_0}
    \end{equation}
    \section{Cross-Sections}
    \begin{equation}
        \bar{y}=\frac{\sum A_i y_i }{A_i}
    \end{equation}
    \begin{equation}
        I=\int_M y^2 dA
    \end{equation}
    \begin{equation}
        I=\sum I_0 + A_i d_i^2
        \quad | \quad d=\bar{y}-y_i
    \end{equation}
    \begin{equation}
        I_{Rect}=\frac{bh^3}{12}
    \end{equation}
    \begin{equation}
        I_{Circ}=\frac{\pi d^4}{64}
    \end{equation}
    \begin{equation}
        Q=\int \limits_{y_{\tau=0}}^{y} y dA
    \end{equation}
    \begin{equation}
        Q=\sum A_i d_i
    \end{equation}
    \section{Bending}
    \begin{equation}
        \phi=\frac{d\theta}{dx}=\frac{d^2y}{dx^2}
    \end{equation}
    \begin{equation}
        \varepsilon=\phi y
    \end{equation}
    \begin{equation}
        M=EI\phi
    \end{equation}
    \section{Euler Buckling}
    \begin{equation}
        P_{crit}=\frac{n^2 \pi^2 EI}{L^2}
    \end{equation}
    \begin{equation}
        r=\sqrt{\frac{I}{A}}
    \end{equation}
    \begin{equation}
        \sigma_{crit}=\frac{\pi^2 E}{\left(\frac{L}{r}\right)^2}
    \end{equation}
    \begin{equation}
        \Delta_{lat}=\frac{\Delta_0}{1-{\frac{P}{P_{crit}}}}
    \end{equation}
    \section{Trusses}
    \begin{equation}
        A>\frac{2F_{max}}{\sigma_y}
    \end{equation}
    \begin{equation}
        I>\frac{3L^2C_{max}}{\pi^2 E}
    \end{equation}
    \begin{equation}
        r>\frac{L}{200}
    \end{equation}
    \begin{equation}
        w_{wind}=2\text{kPa}
    \end{equation}
    \begin{equation}
        A_{wind}=\sum \frac{L_i}{2}
    \end{equation}
    \begin{equation}
        A_{wind, rail} = h\cdot s
    \end{equation}
    \begin{equation}
        \Delta_{brace} = 0.01L
    \end{equation}
    \begin{equation}
        R_{brace}= 0.02P
    \end{equation}
    \begin{equation}
        P^* \cdot \Delta_B=\sum P^*_i \cdot \frac{P_iL_0}{EA}
    \end{equation}
    \section{Beams}
    \subsection{Stress Resultants}
    \begin{equation}
        \sigma=\frac{My}{I}+\frac{N}{A}
    \end{equation}
    \begin{equation}
        \tau_c = \frac{VQ_{max}}{Ib_c}
    \end{equation}
    \begin{equation}
        \tau_g = \frac{VQ_{g}}{Ib_g}
    \end{equation}
    \begin{equation}
        V(x)=\int w(x) dx + P_0
    \end{equation}
    \begin{equation}
        M(x)=\int V(x) dx + M_0
    \end{equation}
    \begin{equation}
        \phi(x) = \frac{M(x)}{EI}
    \end{equation}
    \begin{equation}
        E_{w, design}=\frac{E_{05}}{1.5}
    \end{equation}
    \begin{equation}
        E_{w, deflection}=E_{50}
    \end{equation}
    \subsection{Deflections/Moment Area Theorems}
    \begin{equation}
        \Delta \theta_A^B=\int \limits_{A}^{B} \phi (x) dx
    \end{equation}
    \begin{equation}
        \delta_{BA}=\int \limits_{A}^{B} x \phi(x)dx=\sum \bar{x}_i\cdot A_{\phi,i}
    \end{equation}
    \subsection{Local Buckling}
    \begin{equation}
        \mu=-\frac{\varepsilon_y}{\varepsilon_x}
    \end{equation}
    \begin{equation}
        E_{apparent}=\frac{E}{1-\mu^2}
    \end{equation}
    \begin{equation}
        \sigma_{1}=4\cdot \frac{\pi^2 E}{12(1-\mu^2)}\cdot\left(\frac{t}{b}\right)^2
    \end{equation}
    \begin{equation}
        \sigma_{2}=0.425\cdot \frac{\pi^2 E}{12(1-\mu^2)}\cdot\left(\frac{t}{b}\right)^2
    \end{equation}
    \begin{equation}
        \sigma_{3}=6\cdot \frac{\pi^2 E}{12(1-\mu^2)}\cdot\left(\frac{t}{b}\right)^2
    \end{equation}
    \begin{equation}
        \tau=5\cdot \frac{\pi^2 E}{12(1-\mu^2)}\cdot\left(\left(\frac{t}{b}\right)^2+\left(\frac{t}{h}\right)^2\right)
    \end{equation}
    \pagebreak
    \section{Concrete}
    \subsection{Properties}
    \begin{equation}
        f_t'=0.33\sqrt{f_c'}
    \end{equation}
    \begin{equation}
        E_c=4730\sqrt{f_c'} \quad | \quad f_c<0.4f_c'
    \end{equation}
    \subsection{Flexure}
    \begin{equation}
        T=A_s \varepsilon E_s
    \end{equation}
    \begin{equation}
        \eta=\frac{E_s}{E_c}
    \end{equation}
    \begin{equation}
        \rho = \frac{A_s}{bd}
    \end{equation}
    \begin{equation}
        K=\sqrt{(\eta \rho)^2 +2 \eta \rho}-\eta \rho
    \end{equation}
    \begin{equation}
        j=(1-\frac{K}{3})
    \end{equation}
    \begin{equation}
        f_s=\frac{M}{A_sjd}
    \end{equation}
    \begin{equation}
        f_c=\frac{K}{1-K}\frac{M}{\eta A_s jd}
    \end{equation}
    \subsection{Shear}
    \begin{equation}
        v=\frac{V}{b_wjd}
    \end{equation}
    \begin{equation}
        V_c + V_s = V_r 
    \end{equation}
    \begin{equation}
        V<0.5V_c+0.6V_s<0.125 b_wjd \cdot f_c'
    \end{equation}
    \begin{equation}
        V_{c, weak}=\frac{230\sqrt{f_c'}}{1000+0.9d}b_wjd
    \end{equation}
    \begin{equation}
        \frac{A_s f_y}{b_ws}\geq 0.06 \sqrt{f_c'}
    \end{equation}
    \begin{equation}
        s\leq\frac{A_sf_y}{0.06\cdot b_w \sqrt{f_c'}}
    \end{equation}
    \begin{equation}
        V_{c, strong}=0.18\sqrt{f_c'}b_wjd
    \end{equation}
    \begin{equation}
        V_s=\frac{A_v f_y jd \text{cot}(35\degree)}{s}
    \end{equation}
    \begin{equation}
        s<600\text{mm}
    \end{equation}
    \begin{equation}
        s<0.63d
    \end{equation}
    \begin{equation}
    s_{safe}=\frac{0.6A_vf_yjd\text{cot} (35\degree)}{V_{applied}-0.5V_c}
    \end{equation}
    \begin{equation} s_{safe}=\frac{0.6A_vf_yjd\text{cot}(35\degree)}{V_{applied}-0.5\cdot 0.18\sqrt{f_c'}b_wjd}
    \end{equation}
    
    


\end{multicols}


\end{document}
